% Consolidated 14-phase-transformations from 350eef03dcd4ca0e99eacbc61d97caa887d8c65a; original TeX fragments preserved.

% === BEGIN 14-phase-transformations-185-194.tex blob c052a91ec36ef4989c783b3541d2f15b7bab88ca ===
% Irodov 用户底本 PDF 页92--94，原书印刷页83--85，原题185--194。
% 第185题底本印Δm=0.70g却简答1.2J；第189题简答mJ应为MJ。
\thermophaseproblem{185}{83}
在温度 $t=100^\circ\mathrm C$ 时，
无质量活塞下的圆筒内装有饱和水蒸气。
缓慢推动活塞使
$\Delta m=0.70\,\mathrm g$ 的蒸汽冷凝。
假设水蒸气是理想气体，
忽略水的液态体积，求对气体所做的功。
\exsolution
温度为 $T=373.15\,\mathrm K$；
饱和冷凝过程等温且近似等压
$p_{\rm sat}=p_0$。
冷凝物的气相物质的量
$\Delta\nu=\Delta m/M$，
其原先占据的气相体积为
$\Delta V_{\rm vap}=\Delta\nu RT/p_0$。
忽略液态体积，活塞压缩体积
$|\Delta V|=\Delta\nu RT/p_0$，
所以外界对气体所做的正功
\[
\boxed{A_{\rm on}=p_0|\Delta V|
=\frac{\Delta m RT}{M}}.
\]
取水摩尔质量
$M=18\times10^{-3}\,\mathrm{kg/mol}$，
题面质量 $\Delta m=7.0\times10^{-4}\,\mathrm{kg}$，
得
\[
\boxed{A_{\rm on}\simeq1.21\times10^2\,\mathrm J=120\,\mathrm J}.
\]
\excheck \textbf{底本数值冲突：}
书后简答印同一公式但给“$1.2\,\mathrm J$”。
使用题面清楚印出的 $0.70\,\mathrm g$，
由气体定律得到约 $120\,\mathrm J$，
相差100倍。未擅自把冷凝质量改为 $0.007\,\mathrm g$。

\thermophaseproblem{186}{84}
体积 $V=6.0\,\mathrm L$ 的容器里装有水和
饱和水蒸气，压强 $40\,\mathrm{atm}$、
温度 $250^\circ\mathrm C$。
气相比体积
$v_v'=50\,\mathrm{L/kg}$，
整个气液系统质量
$m=5.0\,\mathrm{kg}$。
求气相的质量及体积。
\exsolution
设气相质量 $m_v$，液相质量
$m_l=m-m_v$，
液态水的比体积记为
$v_l'\simeq1.0\,\mathrm{L/kg}$。
总容积为
\[
V=m_vv_v'+(m-m_v)v_l'
=m v_l'+m_v(v_v'-v_l').
\]
解得
\[
\boxed{m_v=\frac{V-mv_l'}{v_v'-v_l'}}.
\]
代入以升、千克计的数值：
\[
m_v=\frac{6.0-5.0(1.0)}{50-1.0}
\,\mathrm{kg}
\simeq\boxed{0.020\,\mathrm{kg}=20\,\mathrm g}.
\]
气相体积
\[
\boxed{V_v=m_vv_v'\simeq(0.0204)(50)
\,\mathrm L\simeq1.0\,\mathrm L}.
\]
\excheck 不能把全部 $6\,\mathrm L$
都当作气相容积；液相虽比体积小，
但总液体质量约5 kg，仍占约5 L。

\thermophaseproblem{187}{84}
温度 $100^\circ\mathrm C$ 时，
圆筒活塞下的饱和水蒸气初体积
$V_0=5.0\,\mathrm L$。
在等温条件下压缩至
$V=1.6\,\mathrm L$，
假设饱和蒸汽为理想气体，
求凝结成液体的质量。
\exsolution
饱和蒸汽在相变过程维持同温饱和压
$p_0\simeq1\,\mathrm{atm}$。
压缩减少的气相体积为
$\Delta V=V_0-V$；
忽略液相体积，凝结的摩尔数满足
$p_0\Delta V=\Delta\nu RT$。
因此
\[
\boxed{m_{\rm cond}=
\frac{Mp_0(V_0-V)}{RT}}.
\]
用 $M=0.018\,\mathrm{kg/mol}$、
$\Delta V=3.4\times10^{-3}\,\mathrm{m^3}$、
$T=373.15\,\mathrm K$：
\[
\boxed{m_{\rm cond}\simeq2.0\,\mathrm g}.
\]
\excheck 等温饱和压缩不是一团不凝结气体的
$pV=\text{常数}$ 过程；气相摩尔数随冷凝而减少。

\thermophaseproblem{188}{84}
一批饱和蒸汽最初只处于气相，
现等温将其所占总容积压缩 $n$ 倍。
假设饱和蒸汽比体积是液相比体积的
$N$ 倍，其中 $N>n$。
求末态液相占总容积的比例 $\eta$。
另若末态体积恰在
$p$--$V$ 等温相变平台的中点，求 $\eta$。
\exsolution
设液相比体积为 $v_l$，
气相比体积为 $v_v=Nv_l$，
物质总质量为 $m$。
最初全为饱和气相，
$V_0=mv_v$，
压缩后 $V_f=V_0/n$。
设末态液相质量 $m_l$，
气相质量 $m-m_l$：
\[
V_f=m_lv_l+(m-m_l)v_v
=mv_v-m_l(v_v-v_l).
\]
因此
$m_l=(V_0-V_f)/(v_v-v_l)$。
液相\emph{体积}占比
\[
\eta=\frac{m_lv_l}{V_f}
=\frac{(V_0-V_f)v_l}{V_f(v_v-v_l)}
=\boxed{\frac{n-1}{N-1}}.
\]
若末态体积在相变平台两端
$mv_l$ 与 $mv_v$ 的中点，
$V_f=m(v_l+v_v)/2$，
于是
\[
\boxed{\eta=\frac{mv_l}{2V_f}
=\frac1{N+1}}.
\]
\excheck 本题问的是\emph{液相体积占比}，
不是液态水的质量占比；
两者相差可很大。

\thermophaseproblem{189}{84}
质量 $m=1.00\,\mathrm{kg}$ 的水，
在标准大气压下已沸腾，
随后全部变为饱和水蒸气。
视饱和蒸汽为理想气体，
求系统的熵增与内能增量。
\exsolution
取沸点
$T=373.15\,\mathrm K$、
水的汽化潜热
$q\simeq2.26\times10^6\,\mathrm{J/kg}$。
等温定压相变吸热为
$Q=mq$，因此可逆相变的熵增
\[
\boxed{\Delta S=\frac{mq}{T}
\simeq\frac{2.26\times10^6}{373.15}
\simeq6.1\times10^3\,\mathrm{J/K}
=6.1\,\mathrm{kJ/K}}.
\]
由第一定律，
$\Delta U=Q-A$。
液相体积可忽略，
蒸汽的气体功
$A=p\Delta V\simeq\nu RT=mRT/M$。
于是
\[
\boxed{\Delta U=m\left(q-\frac{RT}{M}\right)}.
\]
以 $M=0.018\,\mathrm{kg/mol}$，
得
\[
\boxed{\Delta U\simeq2.1\times10^6\,\mathrm J
=2.1\,\mathrm{MJ}}.
\]
\excheck \textbf{底本单位勘误：}
本版书后印成“$2.1\,\mathrm{mJ}$”；
按题面1 kg水的汽化潜热和印出的内能公式，
必须是 $2.1\,\mathrm{MJ}$，与小写m的毫焦差 $10^9$ 倍。

\thermophaseproblem{190}{84}
质量 $m=20\,\mathrm g$、
初温 $0^\circ\mathrm C$ 的水，
装在绝热圆筒中，活塞无质量，
活塞面积
$S=410\,\mathrm{cm^2}$，
外界压强为标准大气压。
向水传入热量
$Q=20.0\,\mathrm{kJ}$，
求活塞上升高度。
\exsolution
活塞维持近似等压
$p_0=1\,\mathrm{atm}$。
先将水从 $0^\circ\mathrm C$
加热至 $100^\circ\mathrm C$，
需热
$Q_1=mc\Delta T$，
其中 $c\simeq4.18\times10^3\,\mathrm{J/(kg\,K)}$、
$\Delta T=100\,\mathrm K$。
\[
Q_1=(0.020)(4180)(100)\simeq8.36\,\mathrm{kJ}.
\]
余热
$Q_v=Q-Q_1\simeq11.64\,\mathrm{kJ}$
用于部分汽化。
取汽化潜热
$q\simeq2.26\,\mathrm{MJ/kg}$，
气相质量
$m_v=Q_v/q$。
在沸点 $T_b\simeq373\,\mathrm K$，
蒸汽近似理想气体且液相体积很小，
活塞体积增加
$\Delta V\simeq m_vRT_b/(Mp_0)$。
故
\[
\boxed{h\simeq
\frac{(Q-mc\Delta T)RT_b}{qMp_0S}}
\simeq\boxed{0.21\,\mathrm m\approx0.20\,\mathrm m}.
\]
\excheck 原书书后给约 $20\,\mathrm{cm}$；
采用常用表中汽化潜热为定压焓变时，
得到约21 cm。若把 $q$ 定义为内能汽化差，
还须计入气体膨胀功，
本稿明确了所用潜热约定。

\thermophaseproblem{191}{84}
无质量活塞下的绝热圆筒里有
$m_v=1.0\,\mathrm g$ 饱和水蒸气，
外界标准大气压。
往其中注入温度
$t_0=22^\circ\mathrm C$ 的
$m=1.0\,\mathrm g$ 冷水。
忽略活塞摩擦和筒热容，
求活塞下降时大气压所做的功。
\exsolution
初态蒸汽为 $100^\circ\mathrm C$，
系统仍在外压约1 atm下变化。
需要把新注入的水从
$T_0\simeq295\,\mathrm K$
加热到沸点
$T_b\simeq373\,\mathrm K$，
需热量
$Q_{\rm warm}=mc(T_b-T_0)$。
这部分热由一部分蒸汽冷凝释放，
故冷凝质量
\[
m_{\rm cond}=\frac{mc(T_b-T_0)}q.
\]
忽略液态体积，蒸汽体积减少
$\Delta V\simeq
(m_{\rm cond}/M)RT_b/p_0$。
大气压力在活塞下降过程中做正功
\[
\boxed{A=p_0\Delta V=
\frac{mc(T_b-T_0)RT_b}{qM}}.
\]
代入
$m=0.001\,\mathrm{kg}$、
$c=4180\,\mathrm{J/(kg\,K)}$、
$q\simeq2.26\times10^6\,\mathrm{J/kg}$，
\[
\boxed{A\simeq25\,\mathrm J}.
\]
\excheck 需要把冷水加热到与饱和蒸汽相同温度；
不能误认为整个原有1g蒸汽都冷凝。

\thermophaseproblem{192}{84}
已知球形凸液面上的饱和蒸汽压
比平面液面的饱和蒸汽压高
$\Delta p\ll p$，
并满足 Kelvin 近似公式
$\Delta p\simeq(\rho_v/\rho_l)(2\alpha/r)$。
在温度 $t=27^\circ\mathrm C$ 时，
求水滴直径 $d$，使
$\Delta p/p=\eta=1.0\%$；
假设蒸汽为理想气体。
\exsolution
用饱和蒸汽理想气体关系
$p=\rho_vRT/M$，
将 $\rho_v/p=M/(RT)$ 代入：
\[
\eta=\frac{\Delta p}{p}
\simeq\frac{\rho_v}{\rho_l}
\frac{2\alpha}{rp}
=\frac{2\alpha M}{r\rho_lRT}.
\]
水滴直径 $d=2r$，故
\[
\boxed{d\simeq
\frac{4\alpha M}{\eta\rho_lRT}}.
\]
取 $T=300.15\,\mathrm K$、
$\alpha\simeq0.072\,\mathrm{N/m}$、
$M=0.018\,\mathrm{kg/mol}$、
$\rho_l=1000\,\mathrm{kg/m^3}$，
\[
\boxed{d\simeq2.1\times10^{-7}\,\mathrm m
\simeq0.2\,\mathrm{\mu m}}.
\]
\excheck 公式只在蒸汽压相对增量很小的范围内适用；
更小水滴通常需使用非线性的 Kelvin 方程。

\thermophaseproblem{193}{84}
在 $t=100^\circ\mathrm C$ 的饱和水蒸汽下，
求每秒从 $1\,\mathrm{cm^2}$ 水面
离开的蒸汽分子总质量。
假设撞向水面的气相分子中，
仅有 $\eta=3.6\%$ 被液相留住。
\exsolution
饱和状态净蒸发质量通量为零，
因此离开液面的质量通量等于
气相分子被液体\emph{实际吸收}的通量。
饱和汽的撞击壁面质量通量为
\[
j_{\rm in}=\frac14\rho_v\bar v
=p\sqrt{\frac{M}{2\pi RT}},
\]
这是由 Maxwell 分布对法向速度积分得到的。
只保留其中比例 $\eta$，
净被吸收的质量通量为 $\eta j_{\rm in}$，
在动态平衡下同样等于离开液面的质量通量：
\[
\boxed{j_{\rm out}=
\eta p_0\sqrt{\frac{M}{2\pi RT}}}.
\]
用
$p_0=101325\,\mathrm{Pa}$、
$T=373.15\,\mathrm K$、
$M=0.018\,\mathrm{kg/mol}$，
计算
$j_{\rm out}\simeq3.5\,\mathrm{kg/(m^2\,s)}$。
换算 $1\,\mathrm{cm^2}=10^{-4}\,\mathrm{m^2}$：
\[
\boxed{\dot m\simeq3.5\times10^{-4}\,\mathrm{kg/s}
=0.35\,\mathrm{g/s}}
\quad\text{（每平方厘米水面）}.
\]
\excheck 这是饱和平衡下\emph{双向交换}
中的单向蒸发通量，不代表有同样大小的
净蒸发质量损失。

\thermophaseproblem{194}{85}
在温度 $T=2000\,\mathrm K$ 下，
钨丝在高真空中蒸发，
测得单位面积质量损失率
$\mu=1.2\times10^{-13}\,
\mathrm{g/(s\,cm^2)}$。
求该温度下钨的饱和蒸汽压强。
\exsolution
真空中的分子逸出通量
与饱和蒸汽压强之间的 Hertz--Knudsen
理想动力学近似为
\[
\mu=p_{\rm sat}\sqrt{\frac{M}{2\pi RT}}.
\]
故
\[
\boxed{p_{\rm sat}=
\mu\sqrt{\frac{2\pi RT}{M}}}.
\]
换算质量通量：
\[
1\,\mathrm{g/(cm^2\,s)}
=10\,\mathrm{kg/(m^2\,s)},
\quad
\mu=1.2\times10^{-12}\,\mathrm{kg/(m^2\,s)}.
\]
钨摩尔质量
$M\simeq0.184\,\mathrm{kg/mol}$，
因此
\[
\boxed{p_{\rm sat}\simeq
1.2\times10^{-12}
\sqrt{\frac{2\pi(8.314)(2000)}{0.184}}
\simeq9.0\times10^{-10}\,\mathrm{Pa}
=0.90\,\mathrm{nPa}}.
\]
\excheck 采用理想蒸发系数近似为1；
实际蒸发表面的凝结系数可能带来修正。
% === END 14-phase-transformations-185-194.tex ===

% === BEGIN 14-phase-transformations-195-204.tex blob 2e5750747e354960b16afdf86438eccfdee1f7da ===
% Irodov 第2.6节原题195--204，底本PDF页94--95，印刷页85--86。
\thermophaseproblem{195}{85}
假设水分子间的相互吸引作用突然消失，
其余宏观状态变量保持不变，
那么水施加到容器壁上的压强会增加多少？
\exsolution
范德瓦耳斯形式的状态方程可写为
\[
p=p_{\rm rep}-\frac{a}{V_m^2},
\]
其中 $V_m$ 为摩尔体积，
$p_{\rm rep}=RT/(V_m-b)$ 为排除体积修正的
有效斥力压强，
$a/V_m^2$ 代表分子间吸引的压强修正。
若吸引项消失而体积、温度瞬时保持不变，
墙压增加
\[
\boxed{\Delta p=\frac a{V_m^2}
=a\left(\frac{\rho}{M}\right)^2}.
\]
对密度约
$\rho\simeq1.0\,\mathrm{kg/L}$ 的水，
摩尔质量 $M\simeq18\,\mathrm{g/mol}$，
取水的范德瓦耳斯参数
$a\simeq5.5\,\mathrm{atm\,L^2/mol^2}$，
则
$V_m\simeq0.018\,\mathrm{L/mol}$，
\[
\boxed{\Delta p\simeq1.7\times10^4\,\mathrm{atm}}.
\]
\excheck 这是范德瓦耳斯模型的瞬态分子作用
\emph{形式估算}；
现实水移除吸引作用后通常不可能维持原有液态结构。

\thermophaseproblem{196}{85}
液体密度为 $\rho$，
单位质量汽化潜热为 $q$。
假设汽化所需潜热主要用于克服液体的
内聚压强 $p_i$，
且液体服从范德瓦耳斯模型，
求 $p_i$ 的近似值，并估算水的数值。
\exsolution
取质量为 $m$ 的液体，
汽化到密度远小于液体的蒸汽时，
克服内聚压力的机械能约为
$p_i$ 乘初始液体的相应有效体积
$m/\rho$。
按题设将潜热近似视为内聚作用功，
\[
mq\simeq p_i\frac m\rho.
\]
消去质量：
\[
\boxed{p_i\simeq\rho q}.
\]
对于水，
$\rho\simeq10^3\,\mathrm{kg/m^3}$、
$q\simeq2.26\times10^6\,\mathrm{J/kg}$，
故
\[
\boxed{p_i\simeq2.3\times10^9\,\mathrm{Pa}
\simeq2.2\times10^4\,\mathrm{atm}}.
\]
\excheck 题给将潜热全归于内聚功属于近似假定，
真实汽化焓还包含对外做功和分子内能等贡献。

\thermophaseproblem{197}{85}
证明：一摩尔范德瓦耳斯气体在临界状态满足
\[
V_{m,c}=3b,\qquad
p_c=\frac{a}{27b^2},\qquad
T_c=\frac{8a}{27Rb},
\qquad
p_cV_{m,c}=\frac38RT_c.
\]
\exsolution
一摩尔范德瓦耳斯状态方程：
\[
p(V_m,T)=\frac{RT}{V_m-b}-\frac{a}{V_m^2}.
\]
临界等温线出现水平拐点，故
\[
\left(\frac{\partial p}{\partial V_m}\right)_{T_c}
=-\frac{RT_c}{(V_m-b)^2}+\frac{2a}{V_m^3}=0,
\]
\[
\left(\frac{\partial^2p}{\partial V_m^2}\right)_{T_c}
=\frac{2RT_c}{(V_m-b)^3}-\frac{6a}{V_m^4}=0.
\]
由第一式
$RT_c=2a(V_m-b)^2/V_m^3$，
由第二式
$RT_c=3a(V_m-b)^3/V_m^4$。
两式相等，约去公共因子：
$2V_m=3(V_m-b)$，
所以
\[
\boxed{V_{m,c}=3b}.
\]
代入第一式得
$RT_c=8a/(27b)$，
再代入状态方程得
$p_c=a/(27b^2)$。
于是
\[
\boxed{p_cV_{m,c}
=\frac a{27b^2}(3b)
=\frac a{9b}
=\frac38RT_c}.
\]
\excheck 临界点是等温线的\emph{水平拐点}，
因此必须同时使用一阶、二阶导数为零，
不能仅仅令斜率为零。

\thermophaseproblem{198}{85}
已知二氧化碳的临界温度
$T_c=304\,\mathrm K$、
临界压强 $p_c=73\,\mathrm{atm}$。
求范德瓦耳斯常数 $a,b$。
\exsolution
由上一题临界关系，
$p_c=a/(27b^2)$、
$RT_c=8a/(27b)$。
先由两式联立消去 $a$：
\[
\boxed{b=\frac{RT_c}{8p_c}},
\]
再由 $a=27p_cb^2$：
\[
\boxed{a=\frac{27R^2T_c^2}{64p_c}}.
\]
以
$R=0.082057\,\mathrm{L\,atm/(mol\,K)}$
计算：
\[
\boxed{b\simeq0.0427\,\mathrm{L/mol}},
\qquad
\boxed{a\simeq3.6\,\mathrm{atm\,L^2/mol^2}}.
\]
\excheck $a$ 的量纲为压强乘摩尔体积的平方；
$b$ 的量纲为摩尔体积。

\thermophaseproblem{199}{85}
苯 $\mathrm{C_6H_6}$ 的临界温度
$T_c=562\,\mathrm K$、
临界压强 $p_c=47\,\mathrm{atm}$。
若按范德瓦耳斯模型，
求苯临界态的质量比体积 $v_c$。
\exsolution
由临界恒等式
$p_cV_{m,c}=(3/8)RT_c$，
临界摩尔体积为
$V_{m,c}=3RT_c/(8p_c)$。
质量比体积是单位质量物质的体积，
\[
v_c=\frac{V_{m,c}}M
=\boxed{\frac{3RT_c}{8Mp_c}}.
\]
取苯摩尔质量
$M\simeq78\,\mathrm{g/mol}$、
$R=0.082057\,\mathrm{L\,atm/(mol\,K)}$，
得到
\[
\boxed{v_c\simeq4.7\times10^{-3}\,\mathrm{L/g}
=4.7\,\mathrm{cm^3/g}}.
\]
\excheck 临界\emph{质量比体积}
与摩尔体积不同，必须除以摩尔质量 $M$。

\thermophaseproblem{200}{85}
将范德瓦耳斯状态方程改写为约化变量形式，
其中
$\pi=p/p_c$、
$\upsilon=V_m/V_{m,c}$、
$\tau=T/T_c$。
如果气体压强达到临界压强的12倍，
而体积是临界体积的一半，
求温度与临界温度之比。
\exsolution
原方程
$(p+a/V_m^2)(V_m-b)=RT$。
代入临界关系
$V_{m,c}=3b$、
$p_c=a/(27b^2)$、
$RT_c=8a/(27b)$：
\[
p=p_c\pi,\quad V_m=3b\upsilon,\quad T=T_c\tau.
\]
得
\[
\left[p_c\pi+\frac a{9b^2\upsilon^2}\right]
(3b\upsilon-b)=RT_c\tau.
\]
用 $a/(9b^2)=3p_c$、
$RT_c=8p_cb$，化为
\[
\boxed{\left(\pi+\frac3{\upsilon^2}\right)
(3\upsilon-1)=8\tau}.
\]
给定 $\pi=12$、$\upsilon=1/2$，
\[
8\tau=(12+3/(1/2)^2)(3/2-1)
=(12+12)\frac12=12,
\]
故
\[
\boxed{\tau=\frac32=1.5}.
\]
\excheck $\upsilon=1/2>1/3$，
满足范德瓦耳斯方程 $V_m>b$ 的必要体积约束。

\thermophaseproblem{201}{85}
已知水的范德瓦耳斯常数 $a,b$。
求：(a) 质量
$m=1.00\,\mathrm{kg}$ 的水在仍作为液体时
可占据的最大体积；
(b) 饱和水蒸气压强的最大可能值
（范德瓦耳斯模型临界上界）。
\exsolution
范德瓦耳斯模型的液气共存支路在临界点终止。
气、液态临界摩尔体积为
$V_{m,c}=3b$。
对总物质的量
$\nu=m/M$，临界总体积
\[
\boxed{V_{\max,\ell}=\nu(3b)=\frac{3bm}{M}}.
\]
当温度达到临界值时，
饱和蒸汽压强达到临界压强：
\[
\boxed{p_{\rm sat,max}=p_c=\frac a{27b^2}}.
\]
以水在模型中的常用近似参数
$b\simeq0.030\,\mathrm{L/mol}$、
$a\simeq5.5\,\mathrm{atm\,L^2/mol^2}$，
$M=0.018\,\mathrm{kg/mol}$，
\[
\boxed{V_{\max,\ell}\simeq5.0\,\mathrm L},\qquad
\boxed{p_{\rm sat,max}\simeq2.3\times10^2\,\mathrm{atm}}.
\]
\excheck 这里的“最大液体体积”指
范德瓦耳斯\emph{液相支路延至临界点}，
不是任意温度或压力下液态水体积上限。

\thermophaseproblem{202}{85}
把二氧化碳视为范德瓦耳斯气体，
利用其参数 $a,b$ 求临界温度和临界质量密度。
\exsolution
对每摩尔范德瓦耳斯气体，
\[
T_c=\frac{8a}{27Rb},\qquad V_{m,c}=3b.
\]
临界质量密度为
\[
\boxed{\rho_c=\frac M{V_{m,c}}=\frac M{3b}}.
\]
取二氧化碳
$a\simeq3.59\,\mathrm{atm\,L^2/mol^2}$、
$b\simeq0.0427\,\mathrm{L/mol}$、
$M\simeq44.0\,\mathrm{g/mol}$，
以适配的 $R=0.082057$，
\[
\boxed{T_c\simeq3.0\times10^2\,\mathrm K
=0.30\,\mathrm{kK}},
\]
\[
\boxed{\rho_c\simeq\frac{44}{3(0.0427)}
\,\mathrm{g/L}\simeq0.34\,\mathrm{g/cm^3}}.
\]
\excheck 摩尔质量与 $b$ 均须使用同一物质，
单位换算
$1000\,\mathrm{g/L}=1\,\mathrm{g/cm^3}$。

\thermophaseproblem{203}{85}
室温下盛有一部分液态乙醚的密闭容器，
准备将其加热至临界状态。
乙醚临界温度
$T_c=467\,\mathrm K$、
临界压强
$p_c=35.5\,\mathrm{atm}$、
摩尔质量
$M=74\,\mathrm{g/mol}$。
求室温时液态乙醚所占容器体积的比例。
\exsolution
设容器总体积为 $V_{\rm vessel}$，
乙醚质量为 $m$，
若总平均比容需要等于临界比容，
则
$V_{\rm vessel}=mV_{m,c}/M$。
而室温液态乙醚自身占据体积
$V_\ell=m/\rho_\ell$。
所以所求体积分数
\[
f=\frac{V_\ell}{V_{\rm vessel}}
=\frac M{\rho_\ell V_{m,c}}.
\]
范德瓦耳斯临界关系
$V_{m,c}=3RT_c/(8p_c)$，
所以
\[
\boxed{f=\frac{8Mp_c}{3\rho_\ell RT_c}}.
\]
以室温乙醚密度
$\rho_\ell\simeq0.71\,\mathrm{g/cm^3}
=710\,\mathrm{g/L}$ 计算，
\[
\boxed{f\simeq0.25}.
\]
\excheck 乙醚室温密度\emph{并未在题面中给出}，
因此数值需要外部物性表或教材常用近似；
公式本身只依赖临界关系和 $\rho_\ell$。

\thermophaseproblem{204}{85}
证明：图2.5中连接气液两相共存态1与5的
等温等压水平线，
会把范德瓦耳斯等温曲线夹出的
区域 I 与 II 分为相同面积
（Maxwell 等面积构造）。
\exsolution
设共存温度为 $T$、
共存压强为 $p_s$，
液相端摩尔体积为 $V_\ell$，
气相端摩尔体积为 $V_g$。
对于沿平衡态、可逆等温路径的
Helmholtz 自由能有
\[
\mathrm dF=-p\,\mathrm dV
\quad(T=\text{常数}).
\]
共存两相在相同 $T,p_s$ 下，
其摩尔 Gibbs 自由能相同：
\[
F(V_g)+p_sV_g=F(V_\ell)+p_sV_\ell.
\]
故
\[
F(V_g)-F(V_\ell)
=-p_s(V_g-V_\ell).
\]
另一方面沿范德瓦耳斯等温曲线
形式积分：
\[
F(V_g)-F(V_\ell)
=-\int_{V_\ell}^{V_g}p_{\rm vdW}(V,T)\,\mathrm dV.
\]
将两式相等：
\[
\boxed{\int_{V_\ell}^{V_g}
[p_{\rm vdW}(V,T)-p_s]\,\mathrm dV=0}.
\]
这意味着等温曲线位于水平线之上的
有符号面积与位于其下方的面积模相等，
即图示区域 I 与 II 面积相等。
\excheck 等面积规则是由\emph{共存端点的
自由能相等}导出；
不能把范德瓦耳斯中间不稳定分支
误当作真实液气共存时可以稳定走过的物理过程。
% === END 14-phase-transformations-195-204.tex ===

% === BEGIN 14-phase-transformations-205-214.tex blob cbac23af0e7d509bd0462a5cb3bddc435a9d215b ===
% 原书印刷页85--86；本版PDF页94--95。
% 注意底本在第213题后将第214题误排为“224”，后文又继续215。
\thermophaseproblem{205}{85--86}
标准大气压下，
把水过冷到 $t=-20^\circ\mathrm C$。
随后体系自发转变到相平衡态时，
有多大质量分数 $\eta$ 的水结冰？
若希望过冷水全部结冰，
需要初始过冷至多少温度？
\exsolution
设液态过冷水质量为 $m$，
初温 $t<0^\circ\mathrm C$。
按绝热转变模型，结冰放出潜热
$\eta m q_m$，使尚未冻结的物质
从 $t$ 升至平衡凝固点 $0^\circ\mathrm C$。
以液态水比热 $c\simeq4.18\,\mathrm{kJ/(kg\,K)}$
近似，
热平衡为
\[
\eta m q_m\simeq mc|t|,
\]
故
\[
\boxed{\eta\simeq\frac{c|t|}{q_m}}.
\]
取冰融化潜热
$q_m\simeq334\,\mathrm{kJ/kg}$、
$|t|=20\,\mathrm K$：
\[
\boxed{\eta\simeq\frac{4.18\times20}{334}
\simeq0.25=25\%}.
\]
若过冷足以使全部质量结冰，
临界条件 $\eta=1$ 给
\[
\boxed{t_{\rm full}\simeq-\frac{q_m}{c}
\simeq-80^\circ\mathrm C}.
\]
\excheck $-80^\circ\mathrm C$ 是常比热、
无外部换热模型的\emph{形式阈值}，
不声称现实纯水在常压下
能稳定维持到该过冷程度。

\thermophaseproblem{206}{86}
在 $0^\circ\mathrm C$ 附近，
压强升高 $\Delta p=1.00\,\mathrm{atm}$。
冰的质量比体积比液态水大
$\Delta v=0.091\,\mathrm{cm^3/g}$。
求冰的熔点变化。
\exsolution
沿固液共存线，
Clapeyron 方程
\[
\frac{\mathrm dp}{\mathrm dT}
=\frac{q_m}{T(v_{\rm water}-v_{\rm ice})}
=-\frac{q_m}{T\Delta v}.
\]
在压强变化较小时
\[
\boxed{\Delta T\simeq
-\frac{T\Delta v}{q_m}\Delta p}.
\]
将 $\Delta v=0.091\,\mathrm{cm^3/g}
=9.1\times10^{-5}\,\mathrm{m^3/kg}$、
$T\simeq273.15\,\mathrm K$、
$q_m\simeq3.34\times10^5\,\mathrm{J/kg}$ 代入，
\[
\boxed{\Delta T\simeq-7.5\times10^{-3}\,\mathrm K
=-7.5\,\mathrm{mK}}.
\]
\excheck 因液态水比冰密，压强升高利于液态，
熔点降低。本版纯文本提取将体积单位
$\mathrm{cm^3/g}$ 误作 $\mathrm{cm^2/g}$，
已按物理量纲澄清。

\thermophaseproblem{207}{86}
在标准气压下，
如果压强降低 $\Delta p=3.2\,\mathrm{kPa}$
会使水的沸点降低 $\Delta T=0.9\,\mathrm K$，
求原来饱和水蒸汽的质量比体积。
\exsolution
Clapeyron 方程
$\mathrm dp/\mathrm dT=
q_v/[T(v_v-v_\ell)]$。
在沸点附近
$v_v\gg v_\ell$，
故
$\Delta p/\Delta T
\simeq q_v/(T v_v)$。
取压强、温度变化的模：
\[
\boxed{v_v\simeq\frac{q_v\Delta T}{T\Delta p}}.
\]
用
$q_v\simeq2.26\times10^6\,\mathrm{J/kg}$、
$T\simeq373\,\mathrm K$：
\[
\boxed{v_v\simeq
\frac{2.26\times10^6(0.9)}{373(3.2\times10^3)}
\simeq1.7\,\mathrm{m^3/kg}}.
\]
\excheck 这是单位\emph{质量}蒸汽的体积，
而非一摩尔蒸汽的体积；
结果与 $100^\circ\mathrm C$ 附近饱和蒸汽量级一致。

\thermophaseproblem{208}{86}
假设饱和水蒸气为理想气体，
求温度
$101.1^\circ\mathrm C$ 时的饱和蒸气压。
\exsolution
以 $T_0=373.15\,\mathrm K$
处 $p_0=1\,\mathrm{atm}$ 为参考，
温度增量
$\Delta T=1.1\,\mathrm K$。
在忽略液态比体积和汽化潜热温度依赖的近似下，
\[
\frac{\mathrm d\ln p}{\mathrm dT}
=\frac{q_vM}{RT^2}.
\]
对小 $\Delta T$，用一阶展开：
\[
p(T_0+\Delta T)
\simeq p_0\left[
1+\frac{q_vM}{RT_0^2}\Delta T\right].
\]
代入
$q_v\simeq2.26\times10^6\,\mathrm{J/kg}$、
$M=0.018\,\mathrm{kg/mol}$：
\[
\boxed{p\simeq1.04\,\mathrm{atm}}.
\]
\excheck 本题温差仅1.1 K，
可以用线性近似；
若温差较大，则应对 $T^{-2}$ 积分。

\thermophaseproblem{209}{86}
在 $100^\circ\mathrm C$ 时，
少量水与其饱和蒸汽共存于密闭容器。
温度升高 $\Delta T=1.5\,\mathrm K$，
在水的液相比体积可以忽略、
饱和蒸汽为理想气体的条件下，
求饱和蒸汽质量增加的百分比。
\exsolution
容器总体积固定，液相占容积很小。
饱和蒸汽质量为
$m_v=p_{\rm sat}VM/(RT)$，所以
\[
\frac{\mathrm dm_v}{m_v}
=\frac{\mathrm dp_{\rm sat}}{p_{\rm sat}}
-\frac{\mathrm dT}{T}.
\]
Clausius--Clapeyron 近似给
$\mathrm d\ln p_{\rm sat}/\mathrm dT
=q_vM/(RT^2)$。
因此
\[
\frac{\Delta m_v}{m_v}
\simeq
\left(\frac{q_vM}{RT}-1\right)\frac{\Delta T}{T}.
\]
代入 $T=373.15\,\mathrm K$、
$\Delta T=1.5\,\mathrm K$：
\[
\boxed{\frac{\Delta m_v}{m_v}\simeq0.05
\approx5\%}.
\]
\excheck 必须同时计入
饱和蒸汽压随温度上升和
理想气体密度中 $1/T$ 的减小因子。

\thermophaseproblem{210}{86}
已知饱和蒸汽压强在温度 $T_0$ 时为 $p_0$。
假设汽化比潜热 $q_v$ 与温度无关、
液相比体积远小于蒸汽比体积、
饱和蒸汽为理想气体，
求饱和蒸汽压强函数 $p(T)$，
并说明这些近似的适用范围。
\exsolution
Clapeyron 方程
$\mathrm dp/\mathrm dT
=q_v/[T(v_v-v_\ell)]$。
取 $v_v\gg v_\ell$ 且
$v_v=RT/(Mp)$：
\[
\frac1p\frac{\mathrm dp}{\mathrm dT}
=\frac{q_vM}{RT^2}.
\]
在 $T_0$ 到 $T$ 上积分：
\[
\ln\frac{p(T)}{p_0}
=\frac{q_vM}{R}
\int_{T_0}^{T}\frac{\mathrm dT'}{T'^2}
=\frac{q_vM}{R}
\left(\frac1{T_0}-\frac1T\right).
\]
故
\[
\boxed{p(T)=p_0
\exp\left[\frac{q_vM}{R}
\left(\frac1{T_0}-\frac1T\right)\right]}.
\]
\excheck 近似适用于远低于临界温度的
\emph{较窄温区}；
靠近临界点时蒸汽与液态比体积接近，
潜热随温度显著变动，上式失效。

\thermophaseproblem{211}{86}
处于初始标准状况的冰，
被压缩到 $p=640\,\mathrm{atm}$。
假设冰的熔点下降与压强升高成正比，
水与冰的质量比体积差
$\Delta v=0.09\,\mathrm{cm^3/g}$
（液水比冰小）。
求有多少质量分数的冰熔化。
\exsolution
由于压强升高降低熔点，
在绝热压缩与相平衡近似下，
冰体温度需随新的融化温度降低，
一部分冰熔化并吸收潜热。
由 Clapeyron 近似
\[
|\Delta T|\simeq
\frac{T\Delta v}{q_m}\,\Delta p.
\]
降温释放的冰显热约为
$mc_{\rm ice}|\Delta T|$，
用于熔化质量 $\eta m$ 的潜热：
$\eta m q_m\simeq mc_{\rm ice}|\Delta T|$。
因此
\[
\boxed{\eta\simeq
\frac{c_{\rm ice}T\Delta v\,\Delta p}{q_m^2}}.
\]
取
$c_{\rm ice}\simeq2.1\times10^3\,\mathrm{J/(kg\,K)}$、
$T\simeq273\,\mathrm K$、
$\Delta p\simeq639\,\mathrm{atm}$、
$q_m\simeq3.34\times10^5\,\mathrm{J/kg}$，
\[
\boxed{\eta\simeq0.03=3\%}.
\]
\excheck 这里使用了\emph{压缩近似绝热}
与熔点线性下降假设；
若大量压缩功以热形式内化，需进一步修正能量平衡。

\thermophaseproblem{212}{86}
二氧化碳三相点附近的饱和蒸汽压
满足以10为底的经验关系
$\log_{10}p=a-b/T$，
其中 $p$ 以 atm 为单位。
升华支路：
$a_s=9.05,\ b_s=1.80\,\mathrm{kK}$；
汽化支路：
$a_v=6.78,\ b_v=1.31\,\mathrm{kK}$。
求：(a) 三相点温度和压强；
(b) 三相点附近的升华、汽化、熔化质量潜热。
\exsolution
三相点两支路压强相等：
\[
a_s-\frac{b_s}{T_t}
=a_v-\frac{b_v}{T_t}.
\]
解得
\[
\boxed{T_t=\frac{b_s-b_v}{a_s-a_v}
=\frac{1800-1310}{9.05-6.78}
\simeq216\,\mathrm K}.
\]
任一支路回代：
\[
\boxed{p_t=10^{\,a_s-b_s/T_t}
\simeq5.1\,\mathrm{atm}}.
\]
由理想蒸汽和 Clausius--Clapeyron 近似
$\mathrm d\ln p/\mathrm dT
=L_{\rm mol}/(RT^2)$。
题给为\emph{十进制对数}：
\[
\frac{\mathrm d\ln p}{\mathrm dT}
=(\ln10)\frac b{T^2},
\quad L_{\rm mol}=R(\ln10)b.
\]
故每单位质量潜热
$q=R(\ln10)b/M$，
二氧化碳 $M\simeq0.044\,\mathrm{kg/mol}$：
\[
\boxed{q_s\simeq0.78\,\mathrm{kJ/g}},\qquad
\boxed{q_v\simeq0.57\,\mathrm{kJ/g}}.
\]
由三相点焓变相加，
$q_s=q_m+q_v$，所以
\[
\boxed{q_m=q_s-q_v\simeq0.21\,\mathrm{kJ/g}}.
\]
\excheck 必须把 $\log_{10}$ 转为自然对数，
带上 $\ln10$，否则潜热会被低估约2.3倍。

\thermophaseproblem{213}{86}
质量 $m=1.00\,\mathrm{kg}$ 的水，
由 $t_1=10^\circ\mathrm C$ 加热至
$t_2=100^\circ\mathrm C$，
再于沸点全部汽化。
求系统的熵增。
\exsolution
先将液态水由
$T_1=283.15\,\mathrm K$
加热至
$T_2=373.15\,\mathrm K$，
视比热 $c$ 为常数，
熵增为
$m c\ln(T_2/T_1)$。
随后在温度 $T_2$ 等温汽化，
熵增为
$m q_v/T_2$。
两段相加：
\[
\boxed{\Delta S=
m\left[c\ln\frac{T_2}{T_1}+\frac{q_v}{T_2}\right]}.
\]
取水 $c\simeq4.18\,\mathrm{kJ/(kg\,K)}$、
$q_v\simeq2.26\,\mathrm{MJ/kg}$，
\[
\boxed{\Delta S\simeq7.2\,\mathrm{kJ/K}}.
\]
\excheck 绝对温度必须用 K；
先升温后相变的两段熵贡献不能遗漏其中一项。

\thermophaseproblem{214}{86}
冰初温为 $t_1=0^\circ\mathrm C$，
先全部融化，再把水加热至
$t_2=100^\circ\mathrm C$，最后完全汽化。
求系统\emph{单位质量}的熵增。
\exsolution
熔化发生在
$T_m\simeq273.15\,\mathrm K$，
其单位质量熵增
$q_m/T_m$。
液态水升温至
$T_b\simeq373.15\,\mathrm K$，
熵增
$c\ln(T_b/T_m)$。
在沸点汽化，
熵增
$q_v/T_b$。
总单位质量熵增：
\[
\boxed{\Delta s=
\frac{q_m}{T_m}
+c\ln\frac{T_b}{T_m}
+\frac{q_v}{T_b}}.
\]
用
$q_m\simeq334\,\mathrm{J/g}$、
$c\simeq4.18\,\mathrm{J/(g\,K)}$、
$q_v\simeq2257\,\mathrm{J/g}$，
\[
\boxed{\Delta s\simeq8.6\,\mathrm{J/(g\,K)}}.
\]
\excheck \textbf{原书题号排印：}
在本版正文 PDF 印刷页86，
第213题之后此题被误印为“224”，
下题紧接着却是“215”；书后答案亦对应第214题。
为使题库和第213、215题连续，
本稿按\emph{第214题}录入并保留此处勘误。
% === END 14-phase-transformations-205-214.tex ===

% === BEGIN 14-phase-transformations-215-219.tex blob 4db7dcfc8346ca21209b16f557e4917a5dddc345 ===
% 第2.6节最后五题，原书印刷页86--87，PDF页95--96。
\thermophaseproblem{215}{86}
把一块质量 $m=90\,\mathrm g$、
初温 $t_1=90^\circ\mathrm C$ 的铜块，
放进含有 $50\,\mathrm g$、
温度 $-3^\circ\mathrm C$ 冰的量热器中。
求达到热平衡时\emph{铜块自身}的熵变。
\exsolution
铜块从 $90^\circ\mathrm C$ 冷却释放热量。
其最多冷却至 $0^\circ\mathrm C$ 可放出
$m c_{\rm Cu}(90\,\mathrm K)
\simeq(90)(0.39)(90)\,\mathrm J
\simeq3.2\,\mathrm{kJ}$。
这些热量先把冰从 $-3^\circ\mathrm C$
加热到 $0^\circ\mathrm C$，
耗热约
$(50)(2.1)(3)\simeq315\,\mathrm J$；
余热只能熔化约
$(3200-315)/(334\,\mathrm{J/g})
\simeq8.6\,\mathrm g$ 冰，
少于原有50g，故末态仍有冰，
平衡温度恰为 $T_f=273.15\,\mathrm K$。
铜块的熵差
\[
\boxed{\Delta S_{\rm Cu}=
m c_{\rm Cu}\ln\frac{T_f}{T_1}}.
\]
取 $T_1=363.15\,\mathrm K$：
\[
\boxed{\Delta S_{\rm Cu}\simeq
(90)(0.39)\ln\frac{273.15}{363.15}
\simeq-10\,\mathrm{J/K}}.
\]
\excheck 本题只问铜块熵变，可为负；
铜块放热与冰吸热后的\emph{整个隔热体系}
熵增仍应非负。

\thermophaseproblem{216}{87}
初温 $0^\circ\mathrm C$ 的冰质量
$m_1=100\,\mathrm g$，
投入含质量
$m_2=100\,\mathrm g$ 水的绝热量热器。
水的初温分别为
(a) $t_2=60^\circ\mathrm C$；
(b) $t_2=94^\circ\mathrm C$。
忽略量热器热容，求平衡时系统总熵增。
\exsolution
记冰融点 $T_m=273.15\,\mathrm K$，
融化潜热 $q_m$，水比热 $c_w$。
(a) 热水冷至 $T_m$ 释放
$Q_h=m_2c_w(T_2-T_m)$，
若
$Q_h<m_1q_m$，
则只融化一部分冰，
末温是 $T_m$，
融冰量 $m_{\rm melt}=Q_h/q_m$。
熵增
\[
\boxed{\Delta S_a=
m_2c_w\ln\frac{T_m}{T_2}
+\frac{m_2c_w(T_2-T_m)}{T_m}}.
\]
用常用物性数值算得
$\Delta S_a\simeq9\,\mathrm{J/K}$，
与原书约 $9.2\,\mathrm{J/K}$ 相符。
(b) 若
$m_2c_w(T_2-T_m)>m_1q_m$，
全部冰先融化，剩余热量使混合水升温。
能量守恒：
\[
m_2c_w(T_2-T_f)=m_1q_m+
m_1c_w(T_f-T_m).
\]
所以
\[
\boxed{T_f=\frac{m_1T_m+m_2T_2-m_1q_m/c_w}
{m_1+m_2}}.
\]
总熵增
\[
\boxed{\Delta S_b=
\frac{m_1q_m}{T_m}
+m_1c_w\ln\frac{T_f}{T_m}
+m_2c_w\ln\frac{T_f}{T_2}}.
\]
用 $c_w\simeq4.18\,\mathrm{J/(g\,K)}$、
$q_m\simeq334\,\mathrm{J/g}$，
得到 $T_f\simeq280\,\mathrm K$，
$\Delta S_b\simeq20\,\mathrm{J/K}$。
\excheck 原书简答给约 $18\,\mathrm{J/K}$，
取不同的冰潜热和比热可影响取整；
本稿保留精确方程和所采用的物性常数，
避免在数值上强行匹配。

\thermophaseproblem{217}{87}
质量 $m=5.0\,\mathrm g$ 的熔融铅，
初温正好是铅熔点
$t_2=327^\circ\mathrm C$。
将其倒入装有大量
$0^\circ\mathrm C$ 冰的量热器中。
铅凝固潜热为
$q_l=22.5\,\mathrm{J/g}$，
铅比热
$c_l=0.125\,\mathrm{J/(g\,K)}$。
求达到热平衡时铅—冰系统总熵增。
\exsolution
冰量很大，平衡温度维持
$T_1=273.15\,\mathrm K$。
铅在熔点
$T_2\simeq600.15\,\mathrm K$ 凝固，
随后固态铅冷却至 $T_1$。
释放热量为
\[
Q=m q_l+mc_l(T_2-T_1),
\]
这些热量使一部分冰在 $T_1$ 融化，
冰的熵增是 $Q/T_1$。
铅的熵变含两段：
\[
\Delta S_{\rm Pb}
=-\frac{mq_l}{T_2}
+mc_l\ln\frac{T_1}{T_2}.
\]
故系统总熵增
\[
\boxed{\Delta S=
mq_l\left(\frac1{T_1}-\frac1{T_2}\right)
+mc_l\left(\frac{T_2}{T_1}-1
-\ln\frac{T_2}{T_1}\right)}.
\]
代入给定参数：
\[
\boxed{\Delta S\simeq0.48\,\mathrm{J/K}>0}.
\]
\excheck 使用的是铅的凝固潜热，
与冰的熔化潜热\emph{不同}；
大冰库只起温度 $T_1$ 的热库作用。

\thermophaseproblem{218}{87}
活塞下有水的饱和蒸汽，
通过压缩或膨胀使蒸汽始终保持饱和、
恰好处在凝结临界状态。
假设蒸汽为理想气体且液相比体积可忽略。
求沿此过程水蒸汽的摩尔热容 $C(T)$，
并计算 $t=100^\circ\mathrm C$ 的数值。
\exsolution
对一摩尔水蒸汽，
状态方程
$p_{\rm sat}V=RT$。
沿饱和曲线微分：
\[
p_{\rm sat}\frac{\mathrm dV}{\mathrm dT}
=R-\frac{RT}{p_{\rm sat}}
\frac{\mathrm dp_{\rm sat}}{\mathrm dT}.
\]
Clausius--Clapeyron 近似：
\[
\frac1{p_{\rm sat}}
\frac{\mathrm dp_{\rm sat}}{\mathrm dT}
\simeq\frac{q_vM}{RT^2}.
\]
代入：
\[
p_{\rm sat}\frac{\mathrm dV}{\mathrm dT}
=R-\frac{q_vM}{T}.
\]
由热力学第一定律
$C=C_V+p_{\rm sat}\mathrm dV/\mathrm dT$，
\[
\boxed{C=C_p-\frac{q_vM}{T}}.
\]
取 $T=373.15\,\mathrm K$、
$q_v\simeq2.26\,\mathrm{MJ/kg}$、
$M=0.018\,\mathrm{kg/mol}$、
水蒸汽 $C_p\simeq35\,\mathrm{J/(mol\,K)}$，
\[
\boxed{C\simeq-74\,\mathrm{J/(mol\,K)}}.
\]
\excheck 负值是\emph{饱和蒸汽沿特定相变边界}
的过程热容，不表示气体稳定态的 $C_p$ 为负。

\thermophaseproblem{219}{87}
一摩尔水起初在温度 $T_1$ 下与极微量
饱和蒸汽平衡，随后全部转变为温度
$T_2$ 下的饱和水蒸气。
忽略水的液相比体积，把蒸汽当作理想气体。
求系统熵增。
\exsolution
熵是状态函数。
构造一条假想可逆路径：
先把一摩尔液态水由 $T_1$
加热至 $T_2$，然后在 $T_2$ 下汽化。
若液态水的定压摩尔热容
$C_{p,\ell}$ 在温区内近似为常数，
升温熵增为
\[
\Delta S_{\rm warm}
=C_{p,\ell}\ln\frac{T_2}{T_1}.
\]
在 $T_2$ 时汽化，
一摩尔汽化吸热为 $Mq_v$，
对应熵增
$\Delta S_{\rm vap}=Mq_v/T_2$。
故
\[
\boxed{\Delta S=
C_{p,\ell}\ln\frac{T_2}{T_1}
+\frac{Mq_v}{T_2}}.
\]
\excheck 原书简答写 $C_p$，
本稿明确它在所选路径中是
\emph{液态水的定压摩尔热容}，
而不是第218题的水蒸汽摩尔热容；
如热容随温度变化，则应用
$\int_{T_1}^{T_2}C_{p,\ell}(T)\mathrm dT/T$。
% === END 14-phase-transformations-215-219.tex ===
